\documentclass[usenames,dvipsnames,aspectratio=169,12pt]{beamer}
\usepackage{tikz}
\usetikzlibrary{tikzmark,positioning}
\usepackage[customcolors]{hf-tikz}
\usepackage{natbib}
\usepackage{bibentry}
\usepackage{microtype, times}
\usepackage{amsmath, amssymb, mathpartir, stmaryrd}
\usepackage{wasysym}
\usepackage{textcomp}
\usepackage[normalem]{ulem}
\usepackage{listings}
\usepackage{color}
\usepackage{xparse}

\NewDocumentCommand{\Arr}{m m}{#1 \to #2}
\NewDocumentCommand{\Fun}{m m}{\mathsf{fun}\, #1 \Rightarrow #2}
\NewDocumentCommand{\App}{m m}{#1(#2)}
\NewDocumentCommand{\Nat}{}{\mathsf{Nat}}
\NewDocumentCommand{\Zero}{}{\mathsf{zero}}
\NewDocumentCommand{\Succ}{m}{\mathsf{succ}(#1)}
\NewDocumentCommand{\Vect}{m m}{\mathsf{Vec}(#1, #2)}
\NewDocumentCommand{\Nil}{}{\mathsf{nil}}
\NewDocumentCommand{\Cons}{m m}{#1 {:} #2}
\NewDocumentCommand{\Lang}{m}{\mathsf{#1}}

\usefonttheme[onlymath]{serif}
\definecolor{JadeGreen}{RGB}{0,168,107}
\definecolor{MunsellPurple}{RGB}{159,0,197}
\definecolor{LightGrey}{gray}{0.9}
\definecolor{CobaltBlue}{RGB}{0,71,171}
\definecolor{FireBrick}{RGB}{228,34,23}
\definecolor{Alabaster}{RGB}{250,250,250}

\hfsetfillcolor{blue!15}
\hfsetbordercolor{blue!15}

\setbeamertemplate{navigation symbols}{}
\setbeamertemplate{footline}{%
  \raisebox{7pt}{\makebox[\paperwidth]{\hfill\makebox[10pt]{\scriptsize\insertframenumber}\quad}}}
\setbeamertemplate{headline}{}
\setbeamercolor{example text}{fg=CobaltBlue}
\setbeamertemplate{itemize items}{$\bullet$}

\title{Implementing a (Modal) Dependent Type Theory}
\author{\textbf{Daniel Gratzer}\inst{1}\\%
  Jonathan Sterling\inst{2} \and Lars Birkedal\inst{1}}
\institute{\inst{1} This University \smiley{} \\ \inst{2} Not This University \frownie{}}
\date{May 24 2019}

\begin{document}
\begin{frame}[noframenumbering]
  \titlepage
\end{frame}

\begin{frame}
  \frametitle{Hello!}
  First things first, the title of this talk is a lie.
  \pause

  \begin{itemize}
  \item I don't want to talk about \emph{modal} dependent type theory.
  \item I'd like to talk about the core idea behind this work.
  \end{itemize}
  \pause
  \bigskip

  An alternative title:
  \begin{center}
    \color{Blue}
    \Large
    Implementing type theory
  \end{center}

\end{frame}

\begin{frame}
  \frametitle{Setting the Scene}

  What is type theory? Type theory is a....
  \only<1>{\setbeamercolor{itemize item}{fg=red}}
  \begin{itemize}
  \item well-behaved programming language with a rich type system.
    \only<1>{\setbeamercolor{itemize item}{fg=CobaltBlue}}
    \only<2>{\setbeamercolor{itemize item}{fg=red}}

  \item framework for reasoning about mathematical objects.
    \only<2>{\setbeamercolor{itemize item}{fg=CobaltBlue}}
    \only<3>{\setbeamercolor{itemize item}{fg=red}}

  \item foundation for mathematics.
  \end{itemize}
  \pause\pause\pause
  \bigskip


  Each viewpoint induces requirements on a good type theory.
\end{frame}

\begin{frame}
  \frametitle{The Fundamental Tension}

  We want a type theory which is suitable for

  \begin{center}
    \textcolor{Red}{studying mathematics}
  \end{center}

  but at the same time still serves as setting for

  \begin{center}
    \textcolor{Blue}{programming}
  \end{center}

\end{frame}

\begin{frame}
  \frametitle{A Concrete Issue}
  Set aside the questions of \textcolor{Red}{mathematics} and \textcolor{Blue}{programming} for a
  second.
  \bigskip

  Type theory has functions
  \[
    \App{M\tikzmark{fun}}{N\tikzmark{arg}}
  \]
  \onslide<1>{When is this application well-typed?}
  \begin{tikzpicture}[
    remember picture,
    overlay,
    expl/.style={draw=orange,fill=orange!30,rounded corners,text width=3cm},
    arrow/.style={red!80!black,ultra thick,->,>=latex}
    ]
    \node<2->[expl]
    (funex)
    at (-4,-1cm)
    {must have type $\Arr{A}{B}$};
    \node<3->[expl]
    (argex)
    at (4,-1cm)
    {must have type $C$};
    \node<4->
    at (0, -2.5cm)
    {\textbf{We must also have $A = C$!}};
    \draw<2->[arrow]
    (funex) to ([xshift=-1ex,yshift=-0.5ex]{pic cs:fun});
    \draw<3->[arrow]
    (argex) to ([xshift=0ex,yshift=-0.5ex]{pic cs:arg});
  \end{tikzpicture}
\end{frame}

% \begin{frame}
%   \frametitle{Some Background}
%   I'd like to explain this problem more precisely. First we need some types.
% \end{frame}

% \begin{frame}
%   \frametitle{Functions and Natural Numbers}
%   Type theory has natural numbers:
%   \[
%     \Zero \quad \Succ{N} \quad M_0 + M_1 \quad M_0 * M_1
%   \]
%   \pause

%   Type theory also has functions:
%   \begin{align*}
%     \Lang{id} &\triangleq \Fun{x}{x}\\
%     \Lang{const} &\triangleq \Fun{x}{\Fun{y}{x}}
%   \end{align*}
%   As an example:
%   \[
%     \App{\App{\Lang{const}}{M_0}}{M_1} = M_0
%   \]
% \end{frame}

% \begin{frame}
%   \frametitle{A More Exotic Type}

%   Type theory also has \emph{length-indexed} lists: vectors.
%   \begin{displaymath}
%     \onslide<1> \underbrace{ \onslide<1->
%       M : \Vect{A\tikzmark{type}}{N\tikzmark{length}}
%       \onslide<1> }_{M \text{ has type } \Vect{A}{N}}
%   \end{displaymath}
%   \begin{tikzpicture}[
%     remember picture,
%     overlay,
%     expl/.style={draw=orange,fill=orange!30,rounded corners,text width=3cm},
%     arrow/.style={red!80!black,ultra thick,->,>=latex}
%     ]
%     \node<2->[expl]
%     (colonex)
%     at (1,-2cm)
%     {``is of type''};
%     \node<3->[expl]
%     (typeex)
%     at (1,4cm)
%     {``a list of $A$s''};
%     \node<3->[expl]
%     (lengthex)
%     at (5, 0.75cm)
%     {``of length $N$''};
%     \draw<2->[arrow]
%     (colonex) to ([xshift=-0.7ex,yshift=-0.1ex]{pic cs:colon});
%     \draw<3->[arrow]
%     (typeex.east) to[out=0,in=230] ([xshift=2ex,yshift=-0.5ex]{pic cs:type});
%     \draw<4->[arrow]
%     (lengthex.west) to ([xshift=1ex,yshift=0.5ex]{pic cs:length});
%   \end{tikzpicture}

% \end{frame}

\end{document}
